\section{Conclusiones}
 Tanto el Estudio de Mapeo Sistemático (SMS) como la Revisión Sistemática de Literatura (SLR) constituyen metodologías fundamentales y complementarias en el ámbito de la investigación científica. Aunque ambas comparten una estructura metodológica similar basada en tres fases principales (planificación, conducción y reporte), sus diferencias en términos de alcance, profundidad y objetivos las convierten en herramientas especializadas para propósitos distintos. El SMS se posiciona como una metodología ideal para obtener una visión panorámica de campos de investigación amplios, identificar tendencias emergentes y detectar brechas en el conocimiento, mientras que la SLR se enfoca en proporcionar respuestas específicas y detalladas a preguntas de investigación puntuales mediante un análisis riguroso y exhaustivo de la evidencia disponible.